% ============================================================
% CURRICULUM VITAE AND TRACK RECORD (max. 4 pages)
% ============================================================
% Header: Lepperød — FIND — Part B1

\newpage
\section{Curriculum Vitae and Track Record}

\subsection*{Personal details}

\begin{tabular}{ll}
\textbf{Name:} & Dr.\ Mikkel Elle Lepperød \\
\textbf{Date of birth:} & 14 September 1985 \\
\textbf{Nationality:} & Norwegian \\
\textbf{ORCID:} & 0000-0002-4262-5549 \\
\textbf{URL:} & \url{https://www.simula.no/people/mikkel} \\
\end{tabular}

\subsection*{Education and key qualifications}

\begin{tabular}{lp{12cm}}
2020 & \textbf{Ph.D.}\ Neuroscience, Faculty of Medicine, University of Oslo, Norway \\
2014 & \textbf{M.Sc.}\ Applied Mathematics, Norwegian University of Life Sciences, Norway \\
\end{tabular}

\subsection*{Current and relevant previous positions}

\begin{tabular}{lp{12cm}}
2026-- & \textbf{Senior Research Scientist} (equiv.\ associate professor), Simula Research Laboratory \\
2023-- & \textbf{Adjunct Associate Professor} (20\%), Dept.\ of Physics, University of Oslo \\
2022--2026 & \textbf{Research Scientist} (equiv.\ assistant professor), Simula Research Laboratory \\
2020--2022 & \textbf{Postdoctoral Fellow}, Dept.\ of Biosciences, University of Oslo \\
\end{tabular}

\noindent\textbf{PhD supervisor:} Prof.\ Torkel Hafting and Prof.\ Marianne Fyhn (University of Oslo). \\
\textbf{Postdoctoral mentor:} Self-directed (group leader from 2020).

\noindent\textbf{Career breaks:} Parental leave (5 months in 2017, 5 months in 2020).

\subsection*{Research outputs}

The following ten outputs demonstrate my trajectory from experimental neuroscience to neuro-inspired AI, establishing the foundations for FIND:

\begin{enumerate}[noitemsep,topsep=2pt]
  \item \textbf{Lepperød ME}, Christensen AC, Lensjø KK, Buccino AP, Yu J, Fyhn M, Hafting T. \textit{Optogenetic pacing of medial septum parvalbumin-positive cells disrupts temporal but not spatial firing in grid cells}. Science Advances, 2021.
  \textit{Role: First author, conceived and led the study. Significance: Established causal evidence for temporal coding in the entorhinal cortex, directly motivating the temporal predictive coding component of PAM.}

  \item \textbf{Lepperød ME}, Stöber T, Hafting T, Fyhn M, Kording KP. \textit{Inferring causal connectivity from pairwise recordings and optogenetics}. PLoS Computational Biology, 2023.
  \textit{Role: First author. Significance: Developed quasi-experimental methods for causal inference in neural circuits, grounding FIND's causal reasoning objectives.}

  \item Schøyen VS$^*$, Pettersen MB$^*$, Holzhausen K, Fyhn M, Malthe-Sørenssen A, \textbf{Lepperød ME}. \textit{Coherently remapping toroidal cells but not grid cells are responsible for path integration in virtual agents}. iScience, 2023.
  \textit{Role: Senior/corresponding author, supervised the project. Significance: Showed that AI agents spontaneously develop biologically consistent spatial cells, establishing our lab's cognitive map expertise.}

  \item Pettersen MB, Schøyen VS, Malthe-Sørenssen A, \textbf{Lepperød ME}. \textit{Learning place cells and remapping by decoding the cognitive map}. eLife, 2025.
  \textit{Role: Senior/corresponding author. Significance: Demonstrated context-dependent memory through orthogonal remapping, key to FIND's modular memory design.}

  \item Pettersen MB, Rogge F, \textbf{Lepperød ME}. \textit{Learning place cell representations and context-dependent remapping}. NeurIPS, 2024.
  \textit{Role: Senior author. Significance: Extended cognitive map learning to context-switching, directly informing WP2's compositional reasoning.}

  \item Kvalsund MK, Ellefsen KO, Glette K, Pontes-Filho S, \textbf{Lepperød ME}. \textit{Sensor movement drives emergent attention and scalability in active neural cellular automata}. Neural Networks, 2026.
  \textit{Role: Senior/corresponding author. Significance: Demonstrated NCA-based inter-module communication with active vision---the primary orchestration mechanism for PAM.}

  \item Guichard E, Reimers F, Kvalsund MK, \textbf{Lepperød M}, Nichele S. \textit{EngramNCA: A neural cellular automaton model of memory transfer}. ALife, 2025.
  \textit{Role: Co-supervisor. Significance: Showed NCA-based one-shot memory transfer between modules, directly enabling PAM's distributed memory architecture.}

  \item Guichard E, Reimers F, Kvalsund MK, \textbf{Lepperød M}, Nichele S. \textit{ARC-NCA: Towards developmental solutions to the Abstraction and Reasoning Corpus}. ALife, 2025. Featured on Machine Learning Street Talk podcast.
  \textit{Role: Co-supervisor. Significance: Demonstrated abstract reasoning through developmental self-organization, informing WP2's compositional reasoning.}

  \item Schøyen VS, Beshkov K, Pettersen MB, Hermansen E, Holzhausen K, Malthe-Sørenssen A, Fyhn M, \textbf{Lepperød ME}. \textit{Hexagons all the way down: Grid cells as a conformal isometric map of space}. PLOS Computational Biology, 2025.
  \textit{Role: Senior/corresponding author. Significance: Provided geometric theory of neural representations, informing HDC integration in PAM.}

  \item Danielli K, \textbf{Lepperød ME}. \textit{A bio-inspired minimal model for non-stationary K-armed bandits}. Biological Cybernetics, 2026.
  \textit{Role: Senior/corresponding author. Significance: Demonstrated bio-inspired adaptive learning under non-stationarity, directly relevant to WP3's open-ended adaptation.}
\end{enumerate}

\subsection*{Significant peer recognition}

\begin{itemize}[noitemsep,topsep=2pt]
  \item \textbf{NORA Award for Early Career Investigator} (2025)---National recognition for contributions to NeuroAI
  \item \textbf{Simula ERC Incubator} (2022)---Selected for institutional support program for ERC candidates
  \item \textbf{Wharton Executive Education} (2022)---``High-Potential Leaders: Accelerating Your Impact''
  \item \textbf{Invited speaker}: University of California San Diego (2022), AI Week Oslo Met (2023), Keynote at Young Researchers in NeuroAI (2023), Simula Bio-inspired Computing Seminar (2024)
  \item \textbf{Workshop organizer}: 4-day NeuroAI workshop on Hurtigruten with 20 international speakers (2024); co-organizer of 3+ NeuroAI workshops connecting national and international labs (2023--2024)
  \item \textbf{Reviewer}: Nature Machine Intelligence, PLoS Computational Biology, eNeuro
  \item \textbf{Board member}: NORA Special Interest Group for NeuroAI Norway
  \item \textbf{Public outreach}: Brain Inspired podcast (2024), Machine Learning Street Talk podcast (2025), NRK interview (2021), opinion pieces in digi.no and forskning.no
\end{itemize}

\subsection*{Project leadership and supervision}

\noindent\textbf{Group leader, BioAI group at Simula} (2020--present): 4.5 MNOK/year, managing 3 PhD students, 1 postdoc. Built the group from inception to a productive research unit with publications at NeurIPS, eLife, Science Advances, and Nature Communications.

\noindent\textbf{Acting co-PI, Bio-inspired neural networks for AI application} (2020--2026): 16 MNOK RCN-funded project, University of Oslo. Co-supervised 3 PhDs.

\noindent\textbf{Supervised to completion}: 2 PhD students (V.\ Schøyen, K.\ Holzhausen), 8+ MSc students (including graduates now at Stanford PhD program and UiO PhD program).

\subsection*{Research software and ML systems engineering}

\noindent\textbf{PAM / PAMNCA codebases:} Large-scale ML engineering behind the preliminary results in this proposal: systematic GPU-cluster experimentation ($\sim$15K GPU-hours on A40s), automated result parsing from 150+ log files, reproducible evaluation pipelines, and continual-learning benchmarking on CORe50 using the official online (NI) protocol.

\noindent\textbf{Community research infrastructure:} Co-developed NEST simulator modules within the Human Brain Project; co-developed the Exdir hierarchical data format and the Expipe experiment-management platform for reproducible neuroscience pipelines.

\noindent\textbf{Benchmark contributions:} Evaluation on official public protocols (CORe50-NI, ARC).

\noindent\textbf{Team profile to complement the PI:} The FIND postdoc will be recruited from the ML systems community (target profile: PhD from an ICLR/ICML-publishing lab with large-scale training experience) to complement the PI's neuroscience background and own the WP2--3 scaling infrastructure.

\subsection*{Additional information}

\noindent\textbf{Community leadership}: Lead role establishing NeuroAI as focus field in the Norway--US Ministry of Education collaboration (2022). Developed and taught the NeuroAI course at University of Oslo (2023).
